% !TeX root = /Users/lucasgoncalves/Library/CloudStorage/OneDrive-SociedadeCampineiradeEducaçãoeInstrução/Modelos/modelo LaTeX - cópia/prova.tex
% Substitua os enunciados e ajuste os espaços de resposta conforme necessário.

\questao{1,0}{Com base no código abaixo, responda às questões:}

\begin{codigoquestao}[language=Java]
public class Exemplo {
    public static void main(String[] args) {
        System.out.println("Olá, mundo!");
    }
}
\end{codigoquestao}

\subquestao{a}{0,75}{Escreva aqui a primeira pergunta.}
\linhasresposta{1}

\subquestao{b}{0,25}{Escreva aqui a segunda pergunta.}
\linhasresposta{2}

\questao{1,0}{Escreva aqui uma questão que exige uma resposta mais longa.}
\respostalinhas{25}

\questao{1,0}{Escreva aqui o enunciado da questão.}
\respostalinhacodigo{25}

\questao{1,0}{Escreva aqui o enunciado da questão.}
\caixalivre{Escreva o diagrama aqui}{25}

\questao{1,0}{Escreva aqui o enunciado da questão.}
\caixavf{
  [ \quad ] & O modelo em cascata é iterativo por definição. \\ \hline
  [ \quad ] & O Scrum utiliza sprints de duração fixa.
}

\questao{1,0}{Escreva aqui o enunciado da questão.}
\caixarelacionar{
  ( 1 ) & Requisito Funcional \\
  ( 2 ) & Requisito Não-Funcional \\
  ( 3 ) & Regra de Negócio
}{
  ( \quad ) & Desempenho mínimo de 2 segundos. \\
  ( \quad ) & Emitir relatório mensal de vendas. \\
  ( \quad ) & Desconto de 10\% para pagamentos via PIX. \\
  ( \quad ) & Autenticação em dois fatores obrigatoria.
}

\questao{1,0}{Escreva aqui o enunciado da questão.}
\begin{codigolacuna}
public class ProcessoService {
    public void executar() {
        if (\lacuna[4cm]) {
            System.out.println("Aprovado");
        } else {
            \lacuna[3.5cm];
        }
    }
}
\end{codigolacuna}

\questao{1,0}{Escreva aqui o enunciado da questão.}
\espacoresposta{8mm}

\questao{1,0}{Escreva aqui o enunciado da questão.}
\espacoresposta{8mm}

\questao{1,0}{Escreva aqui o enunciado da questão.}

